\subsection{Documentazione delle strategie di testing}

\subsubsection{Metodo sotto test: \texttt{updateIssueById}}

Questo metodo aggiorna i campi di una issue esistente.
Per prima cosa controlla che l'utente che ha fatto la richiesta abbia i permessi necessari: deve essere un amministratore, l'assegnatario o l'autore dell'issue.
Dopo il controllo, aggiorna solo i campi presenti nel DTO, gestisce l'assegnazione e le etichette, e salva l'issue nel database.

\paragraph{Classi di equivalenza individuate}

Per progettare i test abbiamo individuato le seguenti classi di equivalenza, divise per parametro di input.

\begin{longtable}{|c|p{3.5cm}|p{5cm}|p{5cm}|}
\hline
\textbf{ID} & \textbf{Parametro} & \textbf{Classe valida} & \textbf{Classe non valida} \\
\hline
\endfirsthead
\hline
\textbf{ID} & \textbf{Parametro} & \textbf{Classe valida} & \textbf{Classe non valida} \\
\hline
\endhead
CE-1 &
Ruolo del richiedente &
Amministratore &
Utente senza relazione con l'issue (non \`e autore, assegnatario o amministratore) \\
\hline
CE-2 &
Ruolo del richiedente &
Autore dell'issue &
/ \\
\hline
CE-3 &
Assegnatario (campo DTO) &
Email di un utente membro del progetto &
Email di un utente che non \`e membro del progetto \\
\hline
CE-4 &
Assegnatario (campo DTO) &
Stringa vuota (rimozione assegnatario) &
/ \\
\hline
CE-5 &
Assegnatario (campo DTO) &
Stesso assegnatario gi\`a presente &
/ \\
\hline
CE-6 &
Etichette (campo DTO) &
Lista con una o pi\`u etichette &
Lista vuota \\
\hline
CE-7 &
Campi aggiornabili (titolo, stato, descrizione, priorit\`a, tipo) &
Tutti i campi valorizzati &
/ \\
\hline
\end{longtable}

\paragraph{Copertura delle classi di equivalenza nei test}

Nella tabella qui sotto si vede quali classi di equivalenza copre ciascun test.

\begin{longtable}{|p{6cm}|p{8.5cm}|}
\hline
\textbf{Test} & \textbf{Classi coperte} \\
\hline
\endfirsthead
\hline
\textbf{Test} & \textbf{Classi coperte} \\
\hline
\endhead
\texttt{AdminSetsFields} &
CE-1 (amministratore), CE-6 (lista vuota di etichette), CE-7 (tutti i campi valorizzati) \\
\hline
\texttt{EmptyAssignee} &
CE-2 (autore), CE-4 (stringa vuota come assegnatario) \\
\hline
\texttt{AuthorAssignsIssue} &
CE-2 (autore), CE-3 (assegnatario valido membro del progetto), CE-5 (cambio di assegnatario da uno gi\`a presente) \\
\hline
\texttt{AdminSetsLabels} &
CE-1 (amministratore), CE-5 (stesso assegnatario), CE-6 (lista con etichette) \\
\hline
\texttt{AssigneeNotProjectMember} &
CE-2 (autore), CE-3 classe non valida (assegnatario non membro del progetto) \\
\hline
\texttt{IntruderUpdates} &
CE-1 classe non valida (utente senza permessi) \\
\hline
\end{longtable}

\paragraph{Criteri di copertura strutturale}

I test sono stati pensati per coprire tutti i rami principali del metodo:

\begin{itemize}
    \item \textbf{Controllo dei permessi}: il metodo controlla se l'utente \`e amministratore, assegnatario o autore. Abbiamo coperto sia il caso in cui l'utente ha i permessi (amministratore in \texttt{AdminSetsFields}, autore in \texttt{EmptyAssignee}) sia il caso in cui non li ha (\texttt{IntruderUpdates}, dove ci aspettiamo una \texttt{ForbiddenException}).

    \item \textbf{Gestione dell'assegnatario}: ci sono tre rami possibili:
    \begin{itemize}
        \item assegnatario nullo, quindi non si fa nulla;
        \item stringa vuota, quindi l'assegnatario viene rimosso (coperto da \texttt{EmptyAssignee});
        \item email valida, quindi si controlla che l'utente sia membro del progetto. Questo ramo \`e coperto nel caso positivo da \texttt{AuthorAssignsIssue} e nel caso negativo da \texttt{AssigneeNotProjectMember}.
    \end{itemize}

    \item \textbf{Gestione delle etichette}: il codice entra nel ramo delle etichette solo se la lista non \`e nulla e non \`e vuota. Il test \texttt{AdminSetsLabels} copre il caso con etichette presenti, mentre \texttt{AdminSetsFields} copre il caso con lista vuota.

    \item \textbf{Creazione della notifica}: la notifica di assegnazione parte solo quando l'assegnatario cambia. I test \texttt{AdminSetsFields} e \texttt{AuthorAssignsIssue} verificano che la notifica venga creata, mentre \texttt{AdminSetsLabels} verifica che non venga creata quando l'assegnatario \`e lo stesso di prima.
\end{itemize}

% -----------------------------------------------------------------

\subsubsection{Metodo sotto test: \texttt{creaIssue}}

Questo metodo crea una nuova issue nel database.
Riceve un DTO con i dati dell'issue e un file opzionale da allegare.
Dopo la creazione, salva l'allegato se presente, associa le etichette e notifica i membri del progetto.

\paragraph{Classi di equivalenza individuate}

\begin{longtable}{|c|p{3.5cm}|p{5cm}|p{5cm}|}
\hline
\textbf{ID} & \textbf{Parametro} & \textbf{Classe valida} & \textbf{Classe non valida} \\
\hline
\endfirsthead
\hline
\textbf{ID} & \textbf{Parametro} & \textbf{Classe valida} & \textbf{Classe non valida} \\
\hline
\endhead
CE-8 &
File allegato &
File presente e non vuoto &
Nessun file (null) \\
\hline
CE-9 &
File allegato &
File salvato correttamente &
File il cui salvataggio genera un errore (IOException) \\
\hline
CE-10 &
Etichette (campo DTO) &
Lista con una o pi\`u etichette &
Lista vuota o nulla \\
\hline
\end{longtable}

\paragraph{Copertura delle classi di equivalenza nei test}

\begin{longtable}{|p{6cm}|p{8.5cm}|}
\hline
\textbf{Test} & \textbf{Classi coperte} \\
\hline
\endfirsthead
\hline
\textbf{Test} & \textbf{Classi coperte} \\
\hline
\endhead
\texttt{NoFileOneLabel} &
CE-8 classe non valida (nessun file), CE-10 (lista con una etichetta) \\
\hline
\texttt{WithFileNoLabels} &
CE-8 (file presente), CE-9 (salvataggio corretto), CE-10 classe non valida (nessuna etichetta) \\
\hline
\texttt{WithFileAndError} &
CE-8 (file presente), CE-9 classe non valida (errore nel salvataggio), CE-10 classe non valida (lista vuota) \\
\hline
\end{longtable}

\paragraph{Criteri di copertura strutturale}

I test coprono i rami principali di \texttt{creaIssue}:

\begin{itemize}
    \item \textbf{Gestione del file allegato}: il metodo controlla se il file \`e presente e non vuoto. Il test \texttt{NoFileOneLabel} copre il caso senza file (quindi \texttt{salvaFile} non viene chiamato), \texttt{WithFileNoLabels} copre il caso in cui il file viene salvato correttamente, e \texttt{WithFileAndError} copre il caso in cui il salvataggio fallisce con una \texttt{IOException}, e ci aspettiamo che venga lanciata una \texttt{FileUploadException}.

    \item \textbf{Gestione delle etichette}: il codice itera sulle etichette solo se la lista non \`e nulla e non \`e vuota. Il test \texttt{NoFileOneLabel} verifica che l'etichetta venga trovata o creata e poi associata all'issue. Gli altri due test (\texttt{WithFileNoLabels} e \texttt{WithFileAndError}) verificano che con lista vuota o assente il servizio etichette non venga chiamato.
\end{itemize}

\subsubsection{Tecniche utilizzate}

Per i test abbiamo usato il framework \textbf{Mockito}, che ci permette di isolare la classe sotto test dalle sue dipendenze.
Tutte le repository e i servizi esterni (come \texttt{IssueRepository}, \texttt{UtentiRepository}, \texttt{AllegatiService}, ecc.) sono sostituiti con dei mock attraverso l'annotazione \texttt{@Mock}.
La classe \texttt{IssueService} viene poi creata con questi mock iniettati tramite \texttt{@InjectMocks}.

In questo modo possiamo testare solo la logica del metodo, senza dipendere dal database o da altri servizi.
Per ogni test, configuriamo i mock con \texttt{when(...).thenReturn(...)} per simulare le risposte delle dipendenze, e alla fine verifichiamo il risultato con le asserzioni e controlliamo le interazioni con i mock usando \texttt{verify}.